% =============================================================================
% Clin d'oeil a la litterature
% Cette partie reste separee de l'histoire et apparait a la fin du document.
% =============================================================================

\chapter*{Clin d'œil à la littérature}
\addcontentsline{toc}{chapter}{Clin d'œil à la littérature}
\markboth{Clin d'œil à la littérature}{Clin d'œil à la littérature}

Cette annexe rassemble les principaux personnages et noms utilisés dans
l'histoire de Kael qui font référence à d'autres œuvres littéraires.

Ces références sont avant tout des clins d'œil. Les personnages rencontrés
par Kael et Kamatsu ne sont pas nécessairement des transpositions fidèles de
leurs modèles : certains n'en reprennent qu'un nom, d'autres quelques traits
de caractère, une profession, un talent ou une idée particulière.

% -----------------------------------------------------------------------------
\section*{Kael}

\textbf{Référence d'origine :} Aucune référence unique

\textbf{Œuvre / auteur :} ---

Kael est un personnage original. Il n'est pas directement inspiré d'un
personnage littéraire particulier. Ce sont plutôt les différentes rencontres
de sa jeunesse qui constituent des clins d'œil à des personnages et univers
appréciés, tout en permettant d'expliquer progressivement certaines de ses
aptitudes et aspirations.

% -----------------------------------------------------------------------------
\section*{Kamatsu}

\textbf{Référence d'origine :} Kamatsu et la famille Shinzawai

\textbf{Œuvre :} \textit{La Trilogie de l'Empire}

\textbf{Auteurs :} Raymond E. Feist et Janny Wurts

Kamatsu est un noble tsurani appartenant à la maison Shinzawai. Dans notre
histoire, le prénom Kamatsu est repris directement, tandis que Shinzawa
constitue une adaptation du nom Shinzawai.

Le personnage de notre histoire est cependant très différent : seule la
référence au nom est conservée. Shinzawa correspond à son nom familial
d'origine, avant qu'il ne soit recueilli par la caravane.

% -----------------------------------------------------------------------------
\section*{Bartiméus}

\textbf{Référence d'origine :} Bartiméus

\textbf{Œuvre :} \textit{La Trilogie de Bartiméus}

\textbf{Auteur :} Jonathan Stroud

\referencebreak
Bartiméus est un djinn extrêmement ancien, intelligent et sarcastique. Il est
indépendant, volontiers vantard et possède une très haute opinion de lui-même.

\referencebreak
Dans notre histoire, Bartiméus est évidemment devenu un chat. Il conserve
cependant plusieurs traits évoquant son homonyme : son indépendance, son
assurance, sa tendance à n'en faire qu'à sa tête et cette impression
permanente qu'il considère probablement les humains comme légèrement moins
intelligents que lui.

Nous lui ajoutons également un certain instinct pour juger les personnes qu'il
rencontre. Les membres de la caravane ont appris à tenir compte de son
comportement : lorsqu'il apprécie spontanément quelqu'un, ils considèrent
généralement cela comme un bon signe ; lorsqu'il évite une personne, ils ont
tendance à se montrer plus prudents.

% -----------------------------------------------------------------------------
\section*{Artis}

\textbf{Référence d'origine :} Artis Valpierre

\textbf{Œuvre :} \textit{La Quête d'Ewilan} / univers de Gwendalavir

\textbf{Auteur :} Pierre Bottero

\referencebreak
Artis Valpierre est un Rêveur, membre d'un ordre de guérisseurs capable
notamment d'utiliser l'Art du Rêve pour soigner.

\referencebreak
Dans notre histoire, Artis devient un guérisseur de passage voyageant quelque
temps avec la caravane. Lorsque Bartiméus revient blessé après s'être attaqué
à plus gros que lui, les deux garçons lui demandent son aide. Ses méthodes
intriguent particulièrement Kael et constituent sa première véritable
rencontre avec les soins et la magie curative.

% -----------------------------------------------------------------------------
\section*{Rielle}

\textbf{Référence d'origine :} Rielle Lazuli

\textbf{Œuvre :} \textit{La Loi du Millénaire}

\textbf{Autrice :} Trudi Canavan

\referencebreak
Rielle est une jeune femme dotée notamment d'une importante sensibilité
artistique et d'un talent pour la peinture.

\referencebreak
Notre Rielle reprend essentiellement cette facette de créatrice et peintre.
Kamatsu l'observe travailler et elle finit par l'inviter à approcher, puis lui
transmet quelques notions de dessin, de peinture et surtout une certaine
manière d'observer et de représenter le monde.

Bartiméus devient également, avec un succès très relatif concernant sa
capacité à rester immobile, l'un de ses modèles.

% -----------------------------------------------------------------------------
\section*{Robin Mangil}

\textbf{Référence d'origine :} Robin M'angil

\textbf{Œuvre :} \textit{Tara Duncan}

\textbf{Autrice :} Sophie Audouin-Mamikonian

\referencebreak
Robin M'angil est un demi-elfe, compagnon de Tara, excellent archer et
combattant. Ses origines elfiques lui confèrent notamment des capacités
physiques et sensorielles remarquables.

\referencebreak
Dans notre histoire, Robin Mangil conserve le statut de demi-elfe et surtout
son talent pour l'archerie. Il remarque les dispositions naturelles de Kael
pour le tir à l'arc et accepte progressivement de lui enseigner une véritable
technique.

C'est également lui qui conseille Kael lorsqu'il dépense ses économies pour
acheter son premier véritable arc d'occasion.

% -----------------------------------------------------------------------------
\section*{Thom Merilin}

\textbf{Référence d'origine :} Thom Merrilin

\textbf{Œuvre :} \textit{La Roue du Temps}

\textbf{Auteur :} Robert Jordan

\referencebreak
Thom Merrilin est un ménestrel itinérant, musicien, chanteur, conteur et
jongleur particulièrement cultivé et expérimenté. Derrière cette apparence
d'artiste voyageur se cache également un homme possédant beaucoup plus de
ressources qu'il ne le laisse paraître.

\referencebreak
Notre Thom Merilin reprend principalement son aspect de ménestrel, de chanteur
et de conteur. Il remarque l'intérêt de Kamatsu pour ses chansons et commence
à lui transmettre quelques rudiments de musique, de chant et de narration.

Il lui laisse surtout une idée importante : apprendre les histoires des autres
est précieux, mais il faut également apprendre à raconter les siennes.

% -----------------------------------------------------------------------------
\section*{Caliban Dalsen --- « Cal »}

\textbf{Référence d'origine :} Caliban Dal Salan, dit Cal

\textbf{Œuvre :} \textit{Tara Duncan}

\textbf{Autrice :} Sophie Audouin-Mamikonian

\referencebreak
Caliban Dal Salan est l'un des compagnons importants de Tara. Issu d'une
famille de Voleurs Patentés, il suit lui-même cette formation et se révèle
particulièrement doué. Les Voleurs Patentés sont bien davantage que de simples
voleurs : ils remplissent notamment des fonctions proches de l'espionnage et
sont spécialisés dans l'infiltration.

Cal est agile, malin, courageux et particulièrement doué pour se faufiler et
pénétrer là où il ne devrait normalement pas pouvoir entrer. Il est également
connu pour son humour et son caractère impertinent. Son Familier est un renard
roux nommé Blondin.

\referencebreak
Dans notre histoire, Caliban Dalsen est une adaptation volontaire de son nom.
Comme dans l'œuvre originale, il préfère rapidement être appelé simplement
Cal. Nous avons surtout repris son talent pour la discrétion, l'infiltration
et l'observation.

C'est donc particulièrement cohérent que Kael passe son temps à essayer de le
surprendre sans jamais y parvenir. Cal finit par lui apprendre à utiliser les
bruits environnants, à choisir où poser les pieds, à suivre quelqu'un et à
observer sans donner l'impression de le faire.

% -----------------------------------------------------------------------------
\section*{Ellundril Chariak}

\textbf{Référence d'origine :} Ellundril Chariakin

\textbf{Œuvres :} \textit{La Quête d'Ewilan}, \textit{Les Mondes d'Ewilan}
et \textit{Le Pacte des Marchombres}

\textbf{Auteur :} Pierre Bottero

\referencebreak
Ellundril Chariakin est une Marchombre légendaire, surnommée la chevaucheuse
de brume. Elle représente une forme d'aboutissement de la Voie des
Marchombres : liberté, fluidité, silence, discrétion et dépassement des
limites. Sa démarche est notamment décrite comme extraordinairement fluide et
silencieuse, donnant presque l'impression qu'elle flotte au-dessus du sol.

Elle est également étroitement associée à la poésie marchombre : plusieurs
poèmes et aphorismes qui lui sont attribués apparaissent dans les romans.

\referencebreak
Dans notre histoire, Ellundril Chariak conserve précisément ces deux aspects.
Kael est fasciné par sa manière de se déplacer, sa souplesse et son silence ;
Kamatsu, lui, est davantage marqué par sa sensibilité poétique et par la façon
dont elle contemple et décrit le monde.

% -----------------------------------------------------------------------------
\section*{Durzo}

\textbf{Référence d'origine :} Durzo Blint

\textbf{Œuvre :} \textit{L'Ange de la Nuit}

\textbf{Auteur :} Brent Weeks

\referencebreak
Durzo Blint est un pisse-culotte, un tueur professionnel disposant de
capacités surnaturelles qui le distinguent d'un simple assassin. Extrêmement
compétent, dangereux, discret et mystérieux, il est considéré comme l'un des
meilleurs dans son domaine.

Pendant des siècles, il est également le porteur du ka'kari noir, un puissant
artefact magique qui lui confère notamment la capacité de devenir invisible
et de se fondre dans les ombres. Il devient le maître d'Azoth, qu'il forme à
l'art des pisse-culottes avant que celui-ci ne prenne l'identité de Kylar
Stern.

\referencebreak
Dans notre histoire, Durzo est volontairement beaucoup plus mystérieux.
Contrairement aux autres voyageurs rencontrés par Kael, il ne lui enseigne
pratiquement rien.

Et c'est précisément ce qui le rend important.

Lorsque Kael le rencontre, il a déjà appris auprès de Cal et observé la
fluidité d'Ellundril. Il possède donc suffisamment d'expérience pour
comprendre certaines techniques de discrétion. Pourtant, lorsqu'il tente de
suivre Durzo, toutes ses connaissances deviennent insuffisantes : il ne
parvient même pas à comprendre comment l'homme réussit à disparaître.

Durzo devient ainsi moins un professeur qu'un idéal à atteindre. Il laisse à
Kael cette idée qu'il existe encore un niveau de maîtrise qu'il ne comprend
pas et prépare ainsi naturellement son éventuelle évolution future vers la
voie du Traqueur des ténèbres.